\section{Robustness Checks}

In this online appendix we present various robustness checks to show that our main findings are robust to various choices in sample and variable construction.

\subsection{COVID}

Both California's statewide stay-at-home order, and Los Angeles's ``Safer at Home'' emergency order, were issued on March 19th, 2020. The last in-person CPC meeting to occur before this was held on March 12th, 2020. The next CPC meeting, scheduled for March 26th, 2020, was canceled. The one following that, scheduled for April 9th, 2020, was also canceled in observance of Passover. The first CPC meeting to be held after the COVID lockdown order occurred on April 23rd, 2020 and it was held via teleconference. Meetings continued remotely until March 23rd, 2023, which was the last CPC meeting to be held remotely. In-person CPC meetings resumed on April 20th, 2023. Even though the CPC resumed meeting in-person in April 2023, members of the public could still opt to attend the meetings remotely via Zoom.

We re-estimate the main ordered logit regression using two samples. In the first sample, we exclude all hearings that were held remotely during the COVID lockdown period. In the second sample, we include \emph{only} hearings that were held remotely during the COVID lockdown period. The results are presented in Columns 1 and 2 of Table \ref{tab_ologit_covid}.\footnote{Note that due to the more limited sample size, year fixed effects are no longer included.} The coefficients on consent calendar and on atypicality are robust to either specification. Interestingly, the coefficients on agenda order and public opposition are sensitive to which sample was used. Public opposition is only statistically significant in the remote sample, whereas agenda order is only statistically significant in the in-person sample. We hypothesize that this is because public comments were much more likely to be submitted in writing during the COVID period, whereas in-person comments were more likely during the non-COVID period.\footnote{As a reminder, we do not observe comments which are made in-person, only comments which are submitted in writing.} This hypothesis is confirmed by Figure \ref{fig_letters_time}, which shows that the number of written public comments was elevated during the remote-hearings period. However, because the remote/in-person split does not substantially alter our interpretation of the main findings, we do not include these results in the main paper.

\subsection{Mayoral Terms}

Our data period spans two Los Angeles Mayoral administrations. First, Eric Garcetti was the mayor at the start of the sample period, and continued until December 11, 2022, when Karen Bass was sworn in as mayor. Because mayors appoint CPC members, it is plausible that the behavior of the CPC overall is influenced by the priorities of the mayor in office. We therefore re-estimate our main ordered logit regression separately for meetings that occurred during Garcetti's term and meetings that occurred during Bass's term. 

The results are presented in Columns 3-4 of Table \ref{tab_ologit_covid}. The coefficients on agenda order, consent calendar, and atypicality are robust to the mayoral term. However, the coefficient on public opposition was significant only during Garcetti's term. It is difficult to disentangle this finding from the previous finding using the COVID/non-COVID split, because Garcetti's term overlapped significantly more with the COVID-years than Bass's term (Bass was only in office for three months before hearings went back to in-person). Because the findings with the mayoral split do not qualitatively alter the interpretation of any of our results, we do not include them in the main paper.

\subsection{Atypicality Measures}

In this section, we repeat the ordered logit regressions with different methods for measuring atypicality. The results are shown in Table \ref{tab_atypicality_robustness}. In column 1, the Euclidean distance is used to measure atypicality. In column 2, the cosine distance is used. In column 3, the Manhattan distance (L1-norm) is used. The main result, that atypicality negatively predicts approval, is qualitatively robust to any of the distance measures, as are the other coefficients.

